% =====================================================================
% Table 2 (PROPOSAL v2, 2026-10-17) --- the simplified matrix, block-tagged
% ---------------------------------------------------------------------
% NOT \input anywhere. Proposal only; the live file is untouched.
% ADOPTION: copy this over tables/table2_closure-conditions.tex
%           (label tab:closure-conditions is unchanged, so no \ref breaks).
%
% Changes vs the live file:
%   1. Block tags \textbf{labour market:} / \textbf{financing:} (Table 3c
%      framing), placed ON THEIR OWN LINE. Reason: with the tag inline the
%      first column's labour condition wraps to two lines while the second
%      column's does not, so the two "financing:" labels land at different
%      heights. Tag-on-own-line + indented condition makes every block
%      exactly two lines, so both labels align by construction (verified by
%      measuring the word boxes in the compiled shim: identical yMin).
%   2. Row labels reduced to \textbf{F1}\newline \scriptsize <gloss>
%      (Table 3c convention); the old "Financing: F1" prefix duplicated the
%      corner cell.
%   3. Corner cell names the axis ("Financing closure"), matching 3c.
%   4. Geometry: 2.2 + 6.2 + 6.2 cm with \tabcolsep=4pt, so the condition
%      lines fit on one line each (this is what makes point 1 work).
%   5. The Leontief column's substitution pin stays on the labour line, as
%      before; it is a production-side condition and carries no tag here
%      (Table 3c is where the three-way tagging lives). The caption says so.
%   6. Caption: "most importantyl" -> "most importantly", the stray "LABOR"
%      removed, the L notation added, and "the two closure conditions" ->
%      "the short-run and the financing condition".
% =====================================================================
\begin{table}[htbp]
\centering
\footnotesize
\setlength{\tabcolsep}{4pt}
\setlength{\parindent}{0pt}
\setlength{\parskip}{0pt}
\renewcommand{\arraystretch}{1.3}
\newcommand{\tblk}[1]{\textbf{#1:}}
\newcommand{\tcond}[1]{\par\hspace{1em}#1}
\begin{tabular}{@{}>{\raggedright\arraybackslash}p{2.2cm}
                  >{\raggedright\arraybackslash}p{6.2cm}
                  >{\raggedright\arraybackslash}p{6.2cm}@{}}
\toprule
\textbf{Financing closure}
& \textbf{Standard CGE}\newline \scriptsize one flexible wage, labour mobile, full employment
& \textbf{Standard Leontief}\newline \scriptsize fixed real wage, fixed coefficients, output absorbs \\
\midrule
\textbf{F1}\newline \scriptsize repurposing of consumption
& \tblk{labour market}\tcond{$L_{i} = L_{i}^{\mathrm{cm}}(p, y_{i}, w)$, \quad $\textstyle\sum_{i} L_{i} = \bar L$}\par
  \vspace{6pt}
  \tblk{financing}\tcond{$\textstyle\sum_{i} p_{i} c_{i}^{h} = E^{h}$}
& \tblk{labour market}\tcond{$w/\Pi(p) = 1$, \quad $(\theta, \epsilon, \sigma) \to 0^{+}$}\par
  \vspace{6pt}
  \tblk{financing}\tcond{$\textstyle\sum_{i} p_{i} c_{i}^{h} = E^{h}$} \\
\midrule
\textbf{F2}\newline \scriptsize refinancing via taxes
& \tblk{labour market}\tcond{$L_{i} = L_{i}^{\mathrm{cm}}(p, y_{i}, w)$, \quad $\textstyle\sum_{i} L_{i} = \bar L$}\par
  \vspace{6pt}
  \tblk{financing}\tcond{$T(p) = \textstyle\sum_{i} p_{i} g_{i}$}
& \tblk{labour market}\tcond{$w/\Pi(p) = 1$, \quad $(\theta, \epsilon, \sigma) \to 0^{+}$}\par
  \vspace{6pt}
  \tblk{financing}\tcond{$T(p) = \textstyle\sum_{i} p_{i} g_{i}$} \\
\midrule
\textbf{F3}\newline \scriptsize external / debt finance
& \tblk{labour market}\tcond{$L_{i} = L_{i}^{\mathrm{cm}}(p, y_{i}, w)$, \quad $\textstyle\sum_{i} L_{i} = \bar L$}\par
  \vspace{6pt}
  \tblk{financing}\tcond{$B_{\mathrm{gov}} = \textstyle\sum_{i} p_{i} g_{i}$}
& \tblk{labour market}\tcond{$w/\Pi(p) = 1$, \quad $(\theta, \epsilon, \sigma) \to 0^{+}$}\par
  \vspace{6pt}
  \tblk{financing}\tcond{$B_{\mathrm{gov}} = \textstyle\sum_{i} p_{i} g_{i}$} \\
\bottomrule
\end{tabular}
\caption{A matrix showing the two standard short-run closures crossed
with the three financing closures. Each cell states the short-run condition
that defines the labour market of that economy and the financing condition
that closes its budget; the Leontief column additionally pins the
substitution elasticities. Notation: $L_{i}$ is sectoral employment,
$L_{i}^{\mathrm{cm}}$ the cost-minimizing labour demand and $\bar L$ the
labour endowment; $\Pi(p)$ is the consumption price index, $E^{h}$
household consumption expenditure, $T(p)$ the programme tax and
$B_{\mathrm{gov}}$ the externally booked programme cost. All symbols are
defined in Section~\ref{sec:kernel} and Appendix~\ref{app:model}.}
\label{tab:closure-conditions}
\end{table}
